\chapter[Trabalhos Relacionados]{Trabalhos Relacionados}
\label{cap:trabalhos-relacionados}

A literatura de verificação de parentesco facial cresceu de forma consistente entre 2010 e 2025. O que era um problema marginal de visão computacional consolidou-se como área com benchmark de referência, competições anuais e contribuições regulares, evoluindo do primeiro classificador baseado em atributos manuais aos modelos com Vision Transformer e atenção cruzada entre faces. Os modelos multimodais aparecem ao final do capítulo como abordagem complementar de avaliação.

\section{Métodos clássicos e datasets fundadores}

O primeiro modelo computacional de verificação visual de parentesco foi o de \citeonline{fang2010towards}, que tratou o problema como classificação binária a partir de atributos faciais extraídos com estrutura pictórica \cite{felzenszwalb2005pictorial}. Os atributos cobriam cores de olhos, pele e cabelo, partes da face, distâncias entre essas partes e histogramas de gradiente, agrupados em um vetor de atributos manual. Os autores montaram um conjunto próprio de 150 pares de pais e filhos, treinaram k-NN e SVM com função de base radial, com 68,6\% de acurácia em validação cruzada.

Em paralelo, surgiram bancos públicos que tornaram a tarefa comparável entre laboratórios. \citeonline{lu2014kinface}
liberaram KinFaceW-I e KinFaceW-II, com cerca de mil pares por dataset distribuídos em quatro tipos de relação. Ambos seguem em uso até hoje para prototipagem rápida ou validação cruzada de cinco dobras. Pouco depois, \citeonline{robinson2016fiw}
introduziram a primeira versão do Families in the Wild, que viraria, dois anos depois, o principal benchmark da área.

\section{Aprendizado profundo e redes siamesas}

Apesar de o reconhecimento facial já operar em deep learning desde 2012, a verificação de parentesco usou pouco essas redes até 2018, sobretudo pela escassez de dados etiquetados. \citeonline{robinson2018visual} mudou esse panorama com a versão expandida do Families in the Wild, com mais de 13.000 imagens distribuídas em mais de mil famílias. Como solução, propôs o uso de redes siamesas com pesos compartilhados, alimentadas por backbones pré-treinados em reconhecimento facial. As escolhas iniciais foram VGG-Face \cite{parkhi2015deep} e ResNet \cite{wen2016discriminative}, removendo as últimas camadas e treinando uma cabeça de classificação sobre a representação resultante. A partir desse trabalho consolidou-se o RFIW Challenge,
competição anual que padronizou splits, métricas e o protocolo Track-I, no qual famílias inteiras são separadas entre treino e teste para evitar vazamento. A maioria dos modelos publicados desde então segue esse protocolo, e foi também o adotado neste trabalho.

\section{Metric learning e perdas contrastivas}

Backbones melhores não bastam, e parte significativa dos ganhos da década seguinte veio de funções de perda mais discriminativas. Em reconhecimento facial, a sequência SphereFace \cite{liu2017sphereface},
CosFace \cite{wang2018cosface}
e ArcFace \cite{deng2019arcface}
reformulou o softmax para impor margens angulares no espaço de embeddings, e essas perdas migraram naturalmente para parentesco. Em paralelo, abordagens com triplet loss \cite{schroff2015facenet}
e Supervised Contrastive \cite{khosla2020supervised}
passaram a ser usadas para aproximar pares de parentes próximos no espaço de embeddings, afastando pares sem vínculo biológico.

Trabalhos posteriores afinaram a perda para o domínio. \citeonline{lyu2023kfc}
propuseram a Fair Contrastive Loss, que adiciona termos para mitigar viés demográfico durante o treino, e relatam cerca de 81\% de acurácia em FIW. Variantes posteriores usam o tipo da relação como sinal extra para estruturar o espaço, em vez de apenas o rótulo binário entre kin e non-kin.

\section{Mecanismos de atenção e Vision Transformers}

A introdução do Vision Transformer \cite{dosovitskiy2020image}
reabriu a questão de qual é o melhor backbone para parentesco facial. Os resultados mais recentes do FIW vêm de modelos que substituem o backbone convolucional pelo ViT pré-treinado. \citeonline{su2023facornet}
propuseram o FaCoR, que adiciona um módulo de cross-attention bidirecional entre os tokens de patch de duas faces, permitindo que o modelo aprenda quais regiões comparar em vez de só comparar embeddings globais. O ganho relatado foi de 4,6\% sobre o melhor resultado anterior. O FaCoR é a inspiração direta do Modelo 02 deste trabalho.

\section{Fusões híbridas e síntese de idade}

Outra linha tenta combinar os indutivos convolucionais do ConvNeXt \cite{liu2022convnet}
com a atenção global do ViT, em modelos de duplo backbone que produzem features locais e globais e as fundem por concatenação, gating ou atenção. Essa estratégia motivou o Modelo 03 deste trabalho, e a literatura recente aponta que aumentar capacidade total de backbone não traz, em geral, ganho proporcional em verificação de parentesco quando os dados de treino são limitados.

Uma terceira linha ataca diretamente a assimetria etária entre pais e filhos. Métodos baseados em síntese de idade, em particular o SAM (Style-based Age Manipulation) \cite{alaluf2022sam},
geram variantes do mesmo rosto em diferentes idades antes da comparação, com a hipótese de que pares com idades parecidas são mais fáceis de classificar. O Modelo 01 deste trabalho explora essa ideia, embora o módulo SAM não tenha sido plenamente ativado nos experimentos por restrições de hardware.

\section{Generalização cross-dataset e fairness}

Apesar dos números altos relatados na literatura, generalização entre datasets continua um problema aberto. Modelos treinados em FIW e testados em KinFaceW costumam perder entre 15\% e 30\% de acurácia \cite{robinson2022survey}, e o problema é ainda mais grave em subgrupos demográficos sub-representados nos conjuntos de treino. Em reconhecimento facial em geral, estudos de fairness apontam diferenças de 5\% a 10\% na taxa de erro entre faces de subgrupos demográficos distintos \cite{buolamwini2018gender}, e por analogia direta espera-se efeito semelhante em verificação de parentesco facial, embora medições específicas no domínio sejam menos comuns.
Esses dois eixos não são atacados frontalmente neste trabalho, mas aparecem como limitação explícita na conclusão.

\section{Modelos multimodais como baselines}

Modelos multimodais generalistas, em particular VLMs treinados em grandes coleções de pares imagem-texto, ganharam tração no fim da década, abrindo a pergunta sobre se essa classe consegue resolver verificação de parentesco em zero-shot, sem treino específico no domínio. O CLIP \cite{radford2021learning}
foi um dos primeiros nesse formato, e modelos posteriores, como GPT-4V e Claude Sonnet, ampliaram o escopo para tarefas mais variadas. Avaliações sistemáticas de VLMs especificamente em verificação de parentesco facial não foram encontradas na literatura consultada para este trabalho, e a maior parte das referências sobre desempenho de VLMs em tarefas faciais cobre identificação ou descrição genérica de atributos. Os resultados parciais reportados em avaliações ad hoc sugerem que esses modelos exploram bem pistas de idade ou gênero, porém falham em distinções genealógicas mais finas como avós e netos.

Por isso, neste trabalho, os VLMs entram apenas como baselines complementares, não como competidores dos modelos supervisionados, e são avaliados em três regimes (zero-shot, adaptação por prompt e prompts direcionados por treinamento) descritos no Capítulo 4.

\section{Posicionamento deste trabalho}

A combinação de famílias presente neste trabalho cobre cinco eixos da literatura citada. O baseline mínimo usa AdaFace pré-treinado em reconhecimento facial com cosseno e threshold em validação, sem qualquer adaptação ao domínio de parentesco, e mede o desempenho de referência da tarefa. O caminho do ViT com cross-attention bidirecional no estilo FaCoR liga-se à linha de \citeonline{su2023facornet} e segue como ponto de referência forte entre os métodos treinados. Como direção exploratória, um DINOv2 ajustado por LoRA com cross-attention diferencial explora adaptação parameter-efficient sobre backbone self-supervised, e um modelo retrieval-augmented mantém uma galeria de pares positivos do treino e recupera vizinhos similares em inferência. A proposta principal combina backbone facial especializado AdaFace, partial unfreeze do último estágio e cinco tokens regionais anatômicos com cross-region attention bidirecional e gate sigmoide, na linha recente de modelos que substituem patches densos por regiões para reduzir o tamanho da matriz de atenção e introduzir prior anatômico. O Modelo 01 com síntese de idade do pré-projeto e o Modelo 03 com fusão híbrida ConvNeXt e ViT ficam registrados como referência histórica e estudo de capacidade, respectivamente.

A diferença em relação à maioria dos trabalhos citados é metodológica. Em vez de propor um único modelo novo e compará-lo a resultados isolados da literatura, este trabalho aplica o mesmo protocolo de avaliação aos cinco modelos supervisionados, com threshold selecionado em validação e aplicado sem reajuste no teste. A avaliação complementar dos modelos multimodais sobre a tarefa multiclasse de classificação de relações é apresentada em separado, sem pretensão de comparação direta às métricas binárias.
